\section{Sequential Pattern Mining in Clinical Pathways}\label{sec:rw_sequences}

While clustering captures who the patients are at a static point in time, it cannot capture the order in which clinical events unfold. The discovery of recurrent ordered event sequences, namely care pathways, is therefore the second analytical pillar of this thesis. The two works reviewed here are representative of the field: one mines and visualizes frequent care pathways directly from EMR data, while the other combines sequence pattern mining with trajectory clustering, an integration whose limitations are precisely what HospitAI aims to overcome.

Perer et al. \cite{perer2015} propose Care Pathway Explorer, a system that mines a set of frequent event sequences from patient EMR data and presents them through interactive visual analytics. The explicit goal is to extract common sequences of medical events, such as diagnoses and treatments, from historical records and to investigate how these sequences correlate with patient outcomes.

Technically, the system adapts a frequent sequence mining algorithm to the realities of EMR data, handling event concurrency, multiple levels of detail, temporal context, and outcome association. The mined patterns are rendered in a dedicated overview-and-flow visual interface, and the approach is validated on a cohort of hyperlipidemic patients with hypertension and diabetes pre-conditions, where the discovered patterns are shown to correspond to clinically meaningful, and in some cases previously unrecognized, knowledge.

Two limitations distinguish this work from the present thesis. First, the discovered patterns are surfaced through a visualization that a trained analyst must explore manually; there is no automated translation of a sequence into a clinical narrative, so interpretation remains a human, visual task. Second, the mining is performed over the cohort as a whole rather than being conditioned on patient sub-phenotypes. In HospitAI, by contrast, sequence mining is scoped per cluster, so that each phenotype yields its own pathway set, and the high-acuity outliers excluded from clustering are mined as a separate scope precisely so that their otherwise rare but high-cost pathways remain visible, as detailed in Chapter~\ref{ch:design}.

Baudrier et al. \cite{baudrier2022} make an important step towards integration by explicitly combining sequence pattern mining with clustering. Their case study identifies homogeneous healthcare-use profiles and treatment sequences in kidney cancer patients on oral anticancer drugs, retrieved from a French medico-administrative database.

The methodology applies a sequence pattern mining algorithm (CP-SPAM) to healthcare-use sequences and then performs hierarchical clustering on principal components over the mined representation, yielding three patient groups that differ significantly in terms of mortality at six and twelve months and in healthcare resource use. This coupling of pattern mining and clustering is conceptually the closest in the literature to the cluster-scoped pathway analytics of HospitAI.

The interpretation step, however, again relies entirely on clinical experts manually examining the resulting clusters and their representative sequences. There is no LLM layer that automatically generates clinical narrative summaries from the mined sequences, and no conversational interface through which a manager could query the analytical output in natural language. In addition, the patterns are not re-ranked by a statistical interestingness measure that suppresses globally frequent but uninformative items; HospitAI addresses this directly through the closed-pattern compression, geometric-mean-lift re-ranking, and Jaccard-based diversity selection formalized in Chapter~\ref{ch:analysis}.
